% eor_angola.tex
\documentclass[11pt,a4paper]{article}
\usepackage[utf8]{inputenc}
\usepackage[T1]{fontenc}
\usepackage[portuguese]{babel}
\usepackage{graphicx}
\usepackage{amsmath,amssymb}
\usepackage{geometry}
\usepackage{hyperref}
\usepackage{caption}
\usepackage{booktabs}
\usepackage{enumitem}
\geometry{margin=2.5cm}

\title{Recuperação Avançada de Petróleo: \\ Estratégia de Planeamento de Projetos nos Campos Petrolíferos Angolanos}
\author{Geraldo Ramos\textsuperscript{a} \and Kyari Yates\textsuperscript{b} \\
\small \textsuperscript{a}Escola de Petrotecnologia e Engenharia, Academia Sonangol, Luanda, Angola. \\
\small \textsuperscript{a}Departamento de Geociências (DGC), ISPTEC, Luanda, Angola. \\
\small \textsuperscript{b}Faculdade de Farmácia e Ciências da Vida, University of Robert Gordon, Aberdeen, UK.}
\date{}

\begin{document}

\maketitle

% Fonte original (PDF convertido): incluir citação do ficheiro fornecido
% Fonte: (PDF convertido)

\begin{abstract}
A maioria dos campos petrolíferos angolanos encontra--se madura ou em maturação. 
Este artigo apresenta uma estratégia abrangente de planeamento para a implementação de técnicas de Recuperação Avançada de Petróleo (EOR) em Angola, desde a triagem e seleção de campo, passando por testes laboratoriais e pilotos, até à implementação e abandono. O documento discute conceitos, classificação de métodos (térmico, químico, miscível/imiscível e outros), requisitos de dados, desafios técnicos e económicos, e propõe a criação de uma Equipa Nacional de Trabalho de EOR (ANEWT) para coordenar iniciativas nacionais. 
\end{abstract}

\noindent \textbf{Palavras-chave:} Recuperação Avançada de Petróleo (EOR), campos maduros, triagem de reservatórios.

\section{Introdução}
A exploração de hidrocarbonetos em Angola começou em 1910 e a primeira produção ocorreu em 1955. Muitos campos estão maduros e existe a necessidade de recuperar petróleo residual e prolongar a vida produtiva dos campos com técnicas EOR. A EOR pode aumentar substancialmente o fator de recuperação e suprir a procura energética, mitigando a dificuldade de descobrir novos campos. A proposta inclui a criação de uma ANEWT na Agência Nacional do Petróleo, Gás e Biocombustíveis (ANPG) para coordenar esforços entre governo, operadores e prestadores de serviços. (Fonte.)

\section{Conceitos e classificação de EOR}
A recuperação de hidrocarbonetos envolve três processos principais: primária, secundária e recuperação avançada (EOR). A EOR refere--se a um conjunto de tecnologias que injectam energia ou fluidos (gás, químicos, calor, micro-organismos) para mobilizar petróleo residual e melhorar o deslocamento dentro dos poros. Os métodos de EOR são geralmente classificados em:

\begin{itemize}[noitemsep]
  \item Térmicos (vapor, combustão in-situ, água quente).
  \item Químicos (polímeros, surfactantes, alcalinos, combinações ASP).
  \item Gás miscível e imiscível (CO\textsubscript{2}, hidrocarbonetos, N\textsubscript{2}).
  \item Outros (microbiano etc.).
\end{itemize}

\subsection{Eficiências de varredura}
A EOR atua sobre eficiência de varredura microscópica (deslocamento ao nível de poro) e macroscópica (varredura do volume do reservatório). O fator de recuperação global pode ser escrito como:
\begin{equation}
RF = PSD \times SE \times D \times T
\end{equation}
onde $PSD$ é o deslocamento em escala de poros, $SE$ a eficiência de varredura, $D$ a drenagem (conectividade) e $T$ o tempo/comercialidade do projecto. (Equação adaptada do artigo.)

\subsection{Número capilar}
O número capilar $N_{ca}$ relaciona forças viscosa e interfacial:
\begin{equation}
N_{ca} = \frac{v \mu}{\sigma \cos\theta}
\end{equation}
em que $v$ é velocidade, $\mu$ viscosidade, $\sigma$ tensão interfacial e $\theta$ ângulo de contacto. A redução da IFT pode diminuir a saturação residual de óleo. (Fonte.)

\section{Oportunidades e desafios da EOR}
\subsection{Oportunidades}
Implementar EOR pode recuperar uma fração substancial do óleo remanescente, prolongar a vida dos campos e gerar sinergias como captura/armazenamento de CO\textsubscript{2}. A cooperação entre governo, NOCs, IOCs e empresas locais cria oportunidades económicas e sociais. (Fonte.)

\subsection{Desafios}
Operações offshore colocam desafios logísticos e de infra-estrutura; testes e implementação exigem dados extensivos (rocha, fluido, PVT, SCAL), instalações de injeção/tratamento e análises económicas para viabilidade. Processos térmicos levantam preocupações ambientais (emissões); químicos implicam custos e problemas de injectividade/adsorção; disponibilidade de CO\textsubscript{2} é desafio para miscibilidade. (Fonte.)

\section{Estratégia de planeamento de EOR}
A implementação de um projecto EOR envolve três fases principais:
\begin{enumerate}
  \item Seleção do projecto EOR (triagem, estudo de campo, simulação física e matemática, estudos de instalações).
  \item Implementação e optimização do projecto (pilotos, monitorização, QA/QC).
  \item Abandono do campo (quando economicamente apropriado), incluindo possibilidade de armazenamento de CO\textsubscript{2}.
\end{enumerate}

\subsection{Gestão de dados e triagem}
A fase de gestão de dados recolhe propriedades de rocha e fluido: profundidade, API, viscosidade, porosidade, permeabilidade, saturações, temperatura, pressão, salinidade, litologia, entre outros. A triagem técnica pode ser feita por métodos convencionais, geológicos ou avançados (data-mining, redes neurais, modelos neuro-fuzzy). (Ver Secção 4 e Tabela de critérios.)

\subsection{Simulação física e matemática}
Testes laboratoriais (SCAL, IFT, molhabilidade, adsorção) e simulações numéricas (história de produção, previsão, análises de sensibilidade) são ambos essenciais para dimensionar e reduzir risco. Recomenda-se a correspondência histórica antes da previsão de performance do EOR.

\subsection{Testes piloto e monitorização}
Os pilotos variam de poço único a padrões múltiplos; devem incluir monitorização rigorosa (pressões, fluxos, análises de amostras, testes de traçadores, monitorização sísmica 4D quando possível) e procedimentos QA/QC.

\section{Implementação e optimização}
Com base nos estudos de viabilidade, desenvolver o Plano de Desenvolvimento de Campo (FDP) que detalhe a sequencia regulatória, infraestrutura de superfície/subsolo e calendário económico. Formação de pessoal e preparação de instalações são críticos, especialmente em Angola onde EOR é ainda emergente.

\section{Abandono do projecto EOR}
O abandono ocorre após esgotamento económico; deve incluir planos de descomissionamento e, quando viável, transição para armazenamento permanente de CO\textsubscript{2}. Preparar condições contratuais e técnicas para reutilização futura do sítio para armazenamento.

\section{Conclusões}
Um plano estratégico para EOR em Angola foi apresentado, salientando a necessidade de dados robustos, triagem adequada, pilotos bem concebidos, formação e coordenação nacional (ANEWT). A cooperação entre governo, operadores e centros de investigação é recomendada para maximizar recuperação e benefícios socioeconómicos.

\section*{Agradecimentos}
Agradecemos à Sonangol EP pela autorização para preparar e publicar este artigo. (Fonte original: artigo convertido).

\begin{thebibliography}{99}
\bibitem{isco2007} ISCO, T. (2007). Recuperação avançada de petróleo (EOR). Teledyne Technologies.
\bibitem{alanazi2007} Al-Anazi, B. D. (2007). Técnicas avançadas de recuperação de petróleo e injeção de nitrogénio. CSEG Recorder.
\bibitem{ramos2018} Ramos, G. A. R. (2018). Triagem baseada em neuro-fuzzy para projetos de EOR. Tese de doutoramento, University of Aberdeen.
\bibitem{lake1989} Lake, L. W. (1989). Recuperação Avançada de Petróleo. Prentice-Hall.
\bibitem{hite2004} Hite, J. R., et al. (2004). Planeamento de projectos EOR. SPE.
% ... (adicione ou ajuste entradas conforme o PDF original)
\end{thebibliography}

\end{document}
